\originalsection{期末作业第3题}
\begin{exercise}\label{final:3}
MIT 18.712 (2010) 期末作业第3题. 设 $\pi:\mathrm{SU}(2)\to\mathrm{SO}(3)$ 为通常的二重覆盖, $A_5\subset\mathrm{SO}(3)$ 为正二十面体的旋转群, $\Gamma=\pi^{-1}(A_5)$. 计算从 $\Gamma$ 的所有循环子群诱导得到的表示的不可约分解. 此群对应仿射 Dynkin 图 $\widetilde E_8$.
\end{exercise}
\begin{solution}
以下解答为编者独立推导的中文解答, 不是官方公布的答案. 我们先构造不可约表示, 再把诱导重数化成循环子群上的特征值重数. 最后的表和公式同时给出每个循环子群、每个一维特征标的答案; 任意循环子群表示的答案再按直和相加.

\textbf{群元素与循环子群.}
记 $z=-I$, 则 $\ker\pi=\{1,z\}$, $|\Gamma|=120$. 任意 $g\in\mathrm{SU}(2)$ 的特征值可写成 $e^{\mathrm i\theta},e^{-\mathrm i\theta}$, 其像是旋转角为 $2\theta$ 的空间旋转. 因而 $A_5$ 中阶数为 $2,3,5$ 的旋转轴在 $\Gamma$ 中分别给出最大循环子群 $H_4,H_6,H_{10}$. 选生成元 $g_m$ 使自然二维表示 $V=\C^2$ 上的特征值为 $\zeta_m,\zeta_m^{-1}$, 其中 $\zeta_m=e^{2\pi\mathrm i/m}$. 取
\[
 H_m=\langle g_m\rangle\quad(m=4,6,10),\qquad
 H_2=\langle z\rangle,\quad H_3=\langle g_6^2\rangle,\quad
 H_5=\langle g_{10}^2\rangle,\quad H_1=\{1\}.
\]
这些给出循环子群的全部共轭类型. 理由是 $A_5$ 的非平凡元素阶数只有 $2,3,5$, 每个旋转轴的旋转组成一个循环群; 而非中心的 $\mathrm{SU}(2)$ 元素有两个不同特征值, 与它交换的矩阵必须保持两条特征直线, 所以其中心化子是该轴对应的圆群. 与 $\Gamma$ 相交恰为上述最大循环群. 同一种轴的子群在 $A_5$ 中共轭, 提升共轭元便得到 $\Gamma$ 中的共轭. 阶数为 $3,5$ 的子群是相应最大循环群中唯一的同阶子群, 阶数为 $2$ 的子群则唯一等于 $\langle z\rangle$. 按阶数 $1,2,3,4,5,6,10$ 排列, 实际子群数分别为 $1,1,10,15,6,10,6$, 共 $49$ 个循环子群.

为计算特征标, 令 $\varphi=(1+\sqrt5)/2$, $\tau=\varphi-1$. 群中共轭类可由 $t=\tr_V(g)$ 区分, 对应的大小如下:
\[
\begin{array}{c|rrrrrrrrr}
t&2&-2&0&1&-1&\varphi&-\varphi&\tau&-\tau\\ \hline
\text{类大小}&1&1&30&20&20&12&12&12&12
\end{array}
\]
这里 $A_5$ 的双换位有 $15$ 个, 三循环有 $20$ 个, 五循环有 $24$ 个. 五循环的中心化子阶数为 $5$, 所以分成两个各有 $12$ 个元素的共轭类, 分别对应旋转角 $2\pi/5$ 与 $4\pi/5$ (同时包括反向旋转). 提升后, 非零迹的两种提升分别形成一个类: 相应中心化子阶数为 $6$ 或 $10$, 类大小为 $20$ 或 $12$. 双换位的两个提升迹都为零, 在 $\Gamma$ 中中心化子阶数为 $4$, 合成一个大小为 $30$ 的类. 这也直接核对了九个类的总大小为 $120$.

\textbf{全部九个不可约表示的构造.}
若 $g$ 在 $V$ 上的特征值为 $\lambda,\lambda^{-1}$, 则它在 $\Sym^nV$ 上的特征值为 $\lambda^{n-2r}$, $0\le r\le n$. 因此其特征标 $P_n(t)$ 满足
\[
 P_0=1,\quad P_1=t,\quad P_{n+1}=tP_n-P_{n-1}.
\]
我们用记号
\[
 \mathbf1=\Sym^0V,\quad A=\Sym^2V,\quad C=\Sym^3V,
 \quad D=\Sym^4V,\quad E=\Sym^5V.
\]
这些实际表示的维数依次是 $1,3,4,5,6$; 另有自然表示 $V$ 的维数为 $2$. 为了取得余下的表示, 注意 $A$ 因 $z$ 的作用为 $1$ 而下降到 $A_5$. 旋转角 $2\theta$ 的三维旋转表示的特征标为 $1+2\cos(2\theta)=t^2-1=P_2(t)$, 因而它的复化与 $A$ 同构. 将这个 $A_5$ 表示与 $S_5$ 中由奇置换共轭给出的自同构复合, 再拉回 $\Gamma$, 得到三维表示 $A'$. 此自同构交换两个五循环类. 例如对 $c=(12345)$, 把 $c$ 变成 $c^2$ 的正规化置换在其五个位置上相当于 $r\mapsto2r$ ($r\in\Z/5\Z$), 在四个非零位置构成四循环, 故为奇置换; 所以奇共轭确实交换两类.

另取 $A_5$ 在五个字母上的置换表示中坐标和为零的四维子空间, 再拉回 $\Gamma$, 记作 $B$. 它的特征标是固定字母数减 $1$: 在恒等元、双换位、三循环、五循环处分别为 $4,0,1,-1$. 最后暂在虚表示环中定义
\[
 V'=V\otimes B-E.
\]
不能仅从这个差式断言 $V'$ 是实际表示; 下面的正交计算会证明这一点.

用递推式展开 $P_n$, 并用 $\varphi^2=\varphi+1$, $\varphi\tau=1$ 简化, 得到完整特征标表:
\begin{center}
\begin{tabular}{c|rrrrrrrrr}
\toprule
 & $2$&$-2$&$0$&$1$&$-1$&$\varphi$&$-\varphi$&$\tau$&$-\tau$\\
\midrule
$\mathbf1$&1&1&1&1&1&1&1&1&1\\
$V$&2&$-2$&0&1&$-1$&$\varphi$&$-\varphi$&$\tau$&$-\tau$\\
$V'$&2&$-2$&0&1&$-1$&$-\tau$&$\tau$&$-\varphi$&$\varphi$\\
$A$&3&3&$-1$&0&0&$\varphi$&$\varphi$&$-\tau$&$-\tau$\\
$A'$&3&3&$-1$&0&0&$-\tau$&$-\tau$&$\varphi$&$\varphi$\\
$B$&4&4&0&1&1&$-1$&$-1$&$-1$&$-1$\\
$C$&4&$-4$&0&$-1$&1&1&$-1$&$-1$&1\\
$D$&5&5&1&$-1$&$-1$&0&0&0&0\\
$E$&6&$-6$&0&0&0&$-1$&1&1&$-1$\\
\bottomrule
\end{tabular}
\end{center}
表头是自然表示的迹 $t$, 不是表示自身的维数或元素阶数. 所有行的数值均为实数. 将类大小记为 $s_t$, 对任意两行 $X,Y$ 计算
\[
 \inner{\chi_X}{\chi_Y}
 =\frac1{120}\sum_t s_t\chi_X(t)\chi_Y(t).
\]
所得九行的 Gram 矩阵就是单位矩阵. 这项计算只涉及表内九项, 也可逐行核对: 例如
\[
 \inner{\chi_V}{\chi_V}
 =\frac{8+40+24(\varphi^2+\tau^2)}{120}=1,
 \qquad \varphi^2+\tau^2=3.
\]
对其余行同样代入 $\varphi+\tau=\sqrt5$, $\varphi-\tau=1$, $\varphi\tau=1$, 便得对角项 $1$、非对角项 $0$. 对实际表示而言, 特征标范数为 $1$ 表明表示不可约. 对虚表示 $V'$, 若其不可约展开为 $\sum_i n_iX_i$, 则 $n_i\in\Z$ 且 $\sum_i n_i^2=1$, 所以只有一项, 系数为 $\pm1$. 又 $\chi_{V'}(1)=2>0$, 故系数为 $1$, $V'$ 确为实际不可约表示. 因而上述九个不可约表示两两不等价. 最后
\[
 1^2+2^2+2^2+3^2+3^2+4^2+4^2+5^2+6^2=120
\]
说明它们已经穷尽全部不可约表示. 特别要注意, 下一阶对称幂不再不可约: 表中逐项相加给出 $\Sym^6V\cong A'\oplus B$.

题面中的 $\widetilde E_8$ 也可由此核对, 无须预先调用图的分类定理. 将每一行乘以 $\chi_V=t$, 再按上表分解, 得
\begin{align*}
V\otimes\mathbf1&=V,& V\otimes V&=\mathbf1\oplus A,& V\otimes A&=V\oplus C,\\
V\otimes C&=A\oplus D,& V\otimes D&=C\oplus E,& V\otimes E&=D\oplus B\oplus A',\\
V\otimes B&=E\oplus V',& V\otimes V'&=B,& V\otimes A'&=E.
\end{align*}
所以 McKay 图是一条 $\mathbf1,V,A,C,D,E,B,V'$ 的链, 并在 $E$ 处接一个端点 $A'$, 即仿射 $E_8$ 图. 这里每条边的含义已经由实际张量积分解证明, 不靠图名来推测重数.

\textbf{三个最大循环子群上的完整权重.}
对 $m=4,6,10$, 定义非负整数 $a_{X,m}(r)$ 为 $g_m$ 在 $X$ 上的特征值 $\zeta_m^r$ 的重数, $0\le r<m$. 下列三个表列出所有这些整数. 对 $\Sym^nV$, 直接把整数 $n,n-2,\ldots,-n$ 按模 $m$ 计数即可. 对 $A'$, 奇共轭在五阶旋转上把角 $2\pi/5$ 换成 $4\pi/5$, 所以 $H_{10}$ 上的三个权重为 $0,4,6$; 在 $H_4,H_6$ 上与 $A$ 相同. 对 $B$, 则用五点置换的循环结构减去一份平凡特征值. 对 $V'$, 从 $B$ 的权重分别加 $1$、减 $1$, 再减去 $E$ 的权重, 得
\[
 a_{V',m}(r)=a_{B,m}(r-1)+a_{B,m}(r+1)-a_{E,m}(r),
\]
其中下标按模 $m$ 理解. 这是从上面的虚表示构造计算特征值重数, 而其为非负整数则由下表实际核实.

\begin{center}
\begin{tabular}{c|rrrr}
\toprule
$X\backslash r$ ($m=4$)&0&1&2&3\\ \midrule
$\mathbf1$&1&0&0&0\\
$V$&0&1&0&1\\
$V'$&0&1&0&1\\
$A$&1&0&2&0\\
$A'$&1&0&2&0\\
$B$&2&0&2&0\\
$C$&0&2&0&2\\
$D$&3&0&2&0\\
$E$&0&3&0&3\\ \bottomrule
\end{tabular}
\qquad
\begin{tabular}{c|rrrrrr}
\toprule
$X\backslash r$ ($m=6$)&0&1&2&3&4&5\\ \midrule
$\mathbf1$&1&0&0&0&0&0\\
$V$&0&1&0&0&0&1\\
$V'$&0&1&0&0&0&1\\
$A$&1&0&1&0&1&0\\
$A'$&1&0&1&0&1&0\\
$B$&2&0&1&0&1&0\\
$C$&0&1&0&2&0&1\\
$D$&1&0&2&0&2&0\\
$E$&0&2&0&2&0&2\\ \bottomrule
\end{tabular}
\end{center}
\begin{center}
\begin{tabular}{c|rrrrrrrrrr}
\toprule
$X\backslash r$ ($m=10$)&0&1&2&3&4&5&6&7&8&9\\ \midrule
$\mathbf1$&1&0&0&0&0&0&0&0&0&0\\
$V$&0&1&0&0&0&0&0&0&0&1\\
$V'$&0&0&0&1&0&0&0&1&0&0\\
$A$&1&0&1&0&0&0&0&0&1&0\\
$A'$&1&0&0&0&1&0&1&0&0&0\\
$B$&0&0&1&0&1&0&1&0&1&0\\
$C$&0&1&0&1&0&0&0&1&0&1\\
$D$&1&0&1&0&1&0&1&0&1&0\\
$E$&0&1&0&1&0&2&0&1&0&1\\ \bottomrule
\end{tabular}
\end{center}
每行之和都是相应表示的维数. 例如 $E=\Sym^5V$ 在 $H_{10}$ 上的权重为 $5,3,1,-1,-3,-5$, 故权重 $5$ 出现两次, 这解释了表中唯一的重数 $2$ 的中间位置. 一般也可从特征标表以离散 Fourier 变换重算每项:
\[
 a_{X,m}(r)=\frac1m\sum_{k=0}^{m-1}
 \chi_X(g_m^k)\zeta_m^{-rk},\qquad
 \tr_V(g_m^k)=2\cos\frac{2\pi k}{m}.
\]
这是因为 $m^{-1}\sum_{k=0}^{m-1}\zeta_m^{(s-r)k}$ 在 $s\equiv r\pmod m$ 时为 $1$, 否则为 $0$.

\textbf{所有诱导表示的答案.}
对 $H_d$ 的生成元 $h_d$, 记一维特征标为 $\eta_{d,j}(h_d)=\zeta_d^j$, $0\le j<d$. 对最大子群取 $h_d=g_d$. 对 $d=2,3,5$, 分别选 $(m,d)=(4,2),(6,3),(10,5)$, 并取 $h_d=g_m^{m/d}$. Frobenius 互反给出
\[
 [\Ind_{H_d}^{\Gamma}\eta_{d,j}:X]
 =[\Res_{H_d}^{\Gamma}X:\eta_{d,j}]
 =\sum_{\substack{0\le r<m\\r\equiv j\ (\mathrm{mod}\ d)}} a_{X,m}(r).
\]
因此对 $d=4,6,10$, 第 $j$ 列就是 $\Ind_{H_d}^{\Gamma}\eta_{d,j}$ 的九个不可约重数; 对 $d=2,3,5$, 将上表中下标模 $d$ 相同的列相加就是答案. 为使较小子群的结果也能直接读取, 列出折叠后的三个表:
\begin{center}
\begin{tabular}{c|rr|rrr|rrrrr}
\toprule
&\multicolumn{2}{c|}{$d=2$}&\multicolumn{3}{c|}{$d=3$}&\multicolumn{5}{c}{$d=5$}\\
$X\backslash j$&0&1&0&1&2&0&1&2&3&4\\ \midrule
$\mathbf1$&1&0&1&0&0&1&0&0&0&0\\
$V$&0&2&0&1&1&0&1&0&0&1\\
$V'$&0&2&0&1&1&0&0&1&1&0\\
$A$&3&0&1&1&1&1&0&1&1&0\\
$A'$&3&0&1&1&1&1&1&0&0&1\\
$B$&4&0&2&1&1&0&1&1&1&1\\
$C$&0&4&2&1&1&0&1&1&1&1\\
$D$&5&0&1&2&2&1&1&1&1&1\\
$E$&0&6&2&2&2&2&1&1&1&1\\ \bottomrule
\end{tabular}
\end{center}
这里每一列始终按行 $\mathbf1,V,V',A,A',B,C,D,E$ 解读为直和重数. 例如
\begin{align*}
\Ind_{H_{10}}^{\Gamma}\eta_{10,0}&=\mathbf1\oplus A\oplus A'\oplus D,\\
\Ind_{H_{10}}^{\Gamma}\eta_{10,5}&=2E,\\
\Ind_{H_5}^{\Gamma}\eta_{5,0}&=\mathbf1\oplus A\oplus A'\oplus D\oplus2E,\\
\Ind_{H_3}^{\Gamma}\eta_{3,1}&=V\oplus V'\oplus A\oplus A'\oplus B\oplus C\oplus2D\oplus2E.
\end{align*}
从平凡子群诱导则是正则表示, 即
\[
 \Ind_{H_1}^{\Gamma}\mathbf1
 =\mathbf1\oplus2V\oplus2V'\oplus3A\oplus3A'
 \oplus4B\oplus4C\oplus5D\oplus6E.
\]
上述七种阶数 $1,2,3,4,5,6,10$ 的子群及其全部特征标都已列全. 若 $U=\bigoplus_{j=0}^{d-1}b_j\eta_{d,j}$ 是任意 $H_d$ 表示, 则 $\Ind U$ 中 $X$ 的重数为上式乘 $b_j$ 后对 $j$ 求和.

最后说明选择生成元与共轭的影响. 将子群共轭, 并同时搬运它的特征标, 不改变诱导表示的同构类型. 若固定一个子群而改取生成元 $\widetilde h=h^u$, $(u,d)=1$, 则以 $\widetilde h$ 标为 $j$ 的特征标, 在旧生成元下标为 $u^{-1}j\pmod d$; 因而只需把相应列置换. 特别地, 表中各列关于 $j\mapsto-j$ 对称, 但对 $H_{10}$, 改取 $g_{10}^3$ 必须按上述规则重新标列, 不能仍把自然表示的权重标成 $\pm1$.

所有表项非负, 每列按不可约维数加权之和分别为 $120/d$, 与诱导表示维数相符. 对阶数 $2$ 的折叠, 无论从 $H_4,H_6$ 还是 $H_{10}$ 进行, 都得到同一张表: $z$ 在 $\mathbf1,A,A',B,D$ 上作用为 $1$, 在 $V,V',C,E$ 上作用为 $-1$. 这些检查同时覆盖中心元、平凡特征标、重复特征值和奇数阶子群的边界情形.
\end{solution}
